% Generated from locale/tr/content/first-order-logic/axiomatic-deduction/rules-and-proofs.tex; do not hand-edit.


\iftag{FOL}
      {\olfileid[tr]{fol}{axd}{rul}}
      {\olfileid[tr]{pl}{axd}{rul}}

\olsection{Kurallar ve \usetoken{P}{derivation}}

\begin{explain}
  Aksiyomatik !!{derivation}s, mantığın belki de en basit !!{derivation}
  sistemi örneğini verir. !!^a{derivation} yalnızca !!{formula}s içeren bir dizidir. Bir dizinin
  !!a{derivation} sayılması için dizideki her !!{formula}, ya bir aksiyom örneği
  olmalı ya da bir veya daha fazla önceki !!{formula}s kullanılarak bir çıkarım
  kuralıyla elde edilmelidir. !!^a{derivation} için !!{derive}s işleminin
  sonucu, dizinin son terimidir; bu terim bir !!{formula} olur.
\end{explain}

\begin{defn}[!!^{derivability}]
$\Gamma$, $\Lang L$ dilinden !!{formula}s içeren bir kümeyse,
$\Gamma$'dan bir \emph{!!{derivation}}, her $i \le n$ için aşağıdaki
koşullardan birini sağlayan sonlu bir $!A_1$, \dots,~$!A_n$ !!{formula}s
dizisidir:
\begin{enumerate}
\item $!A_i \in \Gamma$; veya
\item $!A_i$ bir aksiyomdur; veya
\item $!A_i$, $j < i$ (ve $k < i$) olmak üzere bazı $!A_j$'den (ve
  $!A_k$'den) bir çıkarım kuralıyla elde edilir.
\end{enumerate}
\end{defn}

Neyin doğru bir !!{derivation} sayıldığı, izin verdiğimiz çıkarım kurallarına
(ve elbette aksiyom olarak aldıklarımıza) bağlıdır. Çıkarım kuralı ise belirli
koşullar altında bir !!a{derivation} içindeki $!A_i$ adımının doğru bir çıkarım
adımı olduğunu bildiren bir eğer--ise ifadesidir.

\begin{defn}[Çıkarım kuralı]
Bir \emph{çıkarım kuralı}, $\Gamma$'dan bir !!a{derivation} içinde hangi
adımın doğru bir çıkarım adımı sayılacağına ilişkin yeterli bir koşul verir.
\end{defn}

Örneğin, $!A \in \Gamma$ olan tek elemanlı her $!A$ dizisi açıkça
!!a{derivation} sayıldığından, aşağıdaki çok basit bir çıkarım kuralı olabilir:
\begin{quote}
  $!A \in \Gamma$ ise $!A$, $\Gamma$'dan her !!{derivation} içinde her zaman
  doğru bir çıkarım adımıdır.
\end{quote}
Benzer biçimde $!A$ aksiyomlardan biriyse, tek başına $!A$ da bir
!!a{derivation} oluşturur. Dolayısıyla aşağıdaki de bir çıkarım kuralıdır:
\begin{quote}
  $!A$ bir aksiyomsa doğru bir çıkarım adımıdır.
\end{quote}
Çıkarım kuralı ele alınan adımdan önce gelen !!{formula}s hakkında bilgi kullandığında
durum daha ilginçleşir. Aşağıdaki kurala \emph{modus ponens} denir:
\begin{quote}
  $!B \lif !A$ ile $!B$, !!{derivation} içinde daha önce geçiyorsa $!A$ doğru
  bir çıkarım adımıdır.
\end{quote}
Tek çıkarım kuralı buysa yukarıdaki !!{derivation} tanımımız şuna karşılık
gelir: $!A_1$, \dots,~$!A_n$ dizisi ancak ve ancak bir !!a{derivation} olur;
bunun koşulu, her $i \le n$ için aşağıdakilerden birinin geçerli olmasıdır:
\begin{enumerate}
\item $!A_i \in \Gamma$; veya
\item $!A_i$ bir aksiyomdur; veya
\item bazı $j < i$ için $!A_j$, $!B \lif !A_i$; bazı $k < i$ içinse $!A_k$,
  $!B$'dir.
\end{enumerate}
Son koşul, $!A_i$'nin $!A_j$'den ($!B \lif !A_i$) ve $!A_k$'den ($!B$)
modus ponens yoluyla elde edildiğini söyler. $1$'den $n$'ye ilerleyip her
defasında $\Gamma$'da bulunan, bir aksiyom olan veya bir çıkarım kuralınca
doğru çıkarım adımı olduğu belirlenen !!a{formula}~$!A_i$ bulursak dizinin
tamamı doğru bir !!{derivation} sayılır.

\begin{defn}[!!^{derivability}]
Bir !!{formula}~$!A$, $\Gamma$'dan \emph{!!{derivable}} sayılır; bunun koşulu
$!A$ ile biten ve $\Gamma$'dan başlayan !!a{derivation} bulunmasıdır. Bunu
$\Gamma \Proves !A$ ile gösteririz.
\end{defn}

\begin{defn}[Teoremler]
Bir !!^a{formula}~$!A$, boş kümeden $!A$ için bir !!a{derivation} varsa bir
\emph{teorem}dir. $!A$ teoremse $\Proves !A$, değilse $\Proves/ !A$ yazarız.
\end{defn}

